% ============================================================
\chapter{Proximal Methods}
\label{ch:proximal}
% ============================================================

\section{Introduction}

In the 1960s, Jean-Jacques Moreau, a mathematician at Montpellier, introduced the \emph{proximal operator}: for a convex function $f$, the proximal of $x$ is the point that minimizes $f$ plus a quadratic penalty around $x$. The idea seems innocuous, but it reveals considerable power: where gradient descent requires differentiability, the proximal operator works for any convex lower semicontinuous function, even nonsmooth ones. This observation opened the way to \emph{proximal methods}, which dominate applied optimization today: ISTA and FISTA for the LASSO in statistics, ADMM for distributed optimization, the Douglas--Rachford algorithm for feasibility. Each exploits the decomposition of a complex problem into simple proximal subproblems.

\begin{intuition}
The proximal operator generalizes projection onto a convex set. Where gradient
descent takes a step and then projects, the proximal operator combines both
into a single step: it finds a point close to the argument that also minimizes
the function. It is a trade-off between staying near the current point and
decreasing the objective.
\end{intuition}

\section{Proximal operator}

\begin{definition}[Proximal operator]
\index{proximal operator}
Let $g : \R^n \to \R \cup \{+\infty\}$ be convex, lower semicontinuous, and
proper. The \textbf{proximal operator} of $g$ is:
\[
  \mathrm{prox}_g(v) = \argmin_{x \in \R^n} \left\{
    g(x) + \frac{1}{2} \norm{x - v}^2
  \right\}.
\]
\end{definition}

\begin{proposition}[Existence and uniqueness]
For any convex, l.s.c., proper function $g$, $\mathrm{prox}_g(v)$ exists and
is unique for every $v \in \R^n$.
\end{proposition}

\begin{theorem}[Characterization]
$x^* = \mathrm{prox}_g(v)$ if and only if:
\[
  v - x^* \in \partial g(x^*),
\]
where $\partial g$ denotes the subdifferential of $g$.
\end{theorem}

\begin{proof}
The optimality condition for $\min_x g(x) + \frac{1}{2}\norm{x-v}^2$ is:
$0 \in \partial g(x^*) + (x^* - v)$, i.e., $v - x^* \in \partial g(x^*)$.
\end{proof}

\begin{keyformulas}[title=\textbf{Common proximal operators}]
\begin{itemize}
  \item $g = 0$: $\mathrm{prox}_g(v) = v$.
  \item $g = \iota_C$ (indicator of convex set $C$):
        $\mathrm{prox}_g(v) = \Pi_C(v)$ (projection).
  \item $g = \lambda \norm{\cdot}_1$ (LASSO):
        $[\mathrm{prox}_g(v)]_i = \sign(v_i)
        \max(\abs{v_i} - \lambda, 0)$ (soft thresholding).
  \item $g = \frac{\lambda}{2}\norm{\cdot}^2$:
        $\mathrm{prox}_g(v) = \frac{v}{1 + \lambda}$.
  \item $g = \lambda \norm{\cdot}_2$ (group LASSO):
        $\mathrm{prox}_g(v) = v \max\!\left(1 - \frac{\lambda}{\norm{v}},
        0\right)$.
\end{itemize}
\end{keyformulas}

\begin{definition}[Moreau decomposition]
\index{Moreau decomposition}
\textbf{Moreau's identity} relates the proximal operator of $g$ to that of
its conjugate $g^*$:
\[
  \mathrm{prox}_g(v) + \mathrm{prox}_{g^*}(v) = v.
\]
\end{definition}

\section{Proximal gradient method}

\begin{definition}[Composite problem]
Consider the problem:
\[
  \min_{x \in \R^n} F(x) = f(x) + g(x),
\]
where $f$ is convex differentiable ($\nabla f$ is $L$-Lipschitz) and $g$ is
convex, l.s.c., proper (possibly nondifferentiable).
\end{definition}

\begin{algorithme}[title=\textbf{Proximal Gradient (ISTA)}]
\index{ISTA}
\begin{enumerate}
  \item Choose $x_0$, step size $\gamma \in (0, 1/L]$.
  \item For $k = 0, 1, 2, \ldots$:
  \[
    x_{k+1} = \mathrm{prox}_{\gamma g}\!\left(
      x_k - \gamma \nabla f(x_k)\right).
  \]
\end{enumerate}
\end{algorithme}

\begin{theorem}[ISTA convergence]
With $\gamma = 1/L$, ISTA satisfies:
\[
  F(x_k) - F(x^*) \leq \frac{L \norm{x_0 - x^*}^2}{2k}.
\]
The convergence rate is $\mathcal{O}(1/k)$.
\end{theorem}

\section{Nesterov acceleration: FISTA}

\begin{algorithme}[title=\textbf{FISTA (Fast ISTA)}]
\index{FISTA}
\begin{enumerate}
  \item Choose $x_0 = y_0$, $t_0 = 1$, $\gamma = 1/L$.
  \item For $k = 0, 1, 2, \ldots$:
  \begin{align*}
    x_{k+1} &= \mathrm{prox}_{\gamma g}\!\left(
      y_k - \gamma \nabla f(y_k)\right), \\
    t_{k+1} &= \frac{1 + \sqrt{1 + 4 t_k^2}}{2}, \\
    y_{k+1} &= x_{k+1} + \frac{t_k - 1}{t_{k+1}}(x_{k+1} - x_k).
  \end{align*}
\end{enumerate}
\end{algorithme}

\begin{theorem}[FISTA convergence]
FISTA satisfies:
\[
  F(x_k) - F(x^*) \leq \frac{2L \norm{x_0 - x^*}^2}{(k+1)^2}.
\]
The $\mathcal{O}(1/k^2)$ rate is \textbf{optimal} among first-order methods.
\end{theorem}

\begin{warning}[title=\textbf{FISTA is not monotone}]
Unlike ISTA, the sequence $F(x_k)$ generated by FISTA is not necessarily
nonincreasing. The extrapolation step may temporarily increase the objective.
\end{warning}

\section{ADMM}

\begin{definition}[ADMM]
\index{ADMM}
The \textbf{Alternating Direction Method of Multipliers} solves:
\[
  \min_{x, z} f(x) + g(z) \quad \text{s.t.} \quad Ax + Bz = c.
\]
The iterations are:
\begin{align*}
  x_{k+1} &= \argmin_x \left\{ f(x) + \frac{\rho}{2}
    \norm{Ax + Bz_k - c + u_k}^2 \right\}, \\
  z_{k+1} &= \argmin_z \left\{ g(z) + \frac{\rho}{2}
    \norm{Ax_{k+1} + Bz - c + u_k}^2 \right\}, \\
  u_{k+1} &= u_k + Ax_{k+1} + Bz_{k+1} - c,
\end{align*}
where $u$ is the scaled dual variable and $\rho > 0$.
\end{definition}

\begin{theorem}[ADMM convergence]
Under convexity assumptions on $f$ and $g$ (proper, closed, convex) and
feasibility:
\begin{itemize}
  \item \textbf{Primal residual}: $Ax_k + Bz_k - c \to 0$.
  \item \textbf{Dual residual}: $\rho A^\top B(z_k - z_{k-1}) \to 0$.
  \item \textbf{Objective}: $f(x_k) + g(z_k) \to p^*$.
\end{itemize}
\end{theorem}

\begin{remark}
ADMM is particularly effective when the $x$- and $z$-subproblems admit
closed-form solutions (e.g., least squares + soft thresholding for LASSO).
\end{remark}

\section{Douglas-Rachford splitting}

\begin{definition}[Douglas-Rachford]
\index{Douglas-Rachford}
To solve $\min_x f(x) + g(x)$, the Douglas-Rachford algorithm iterates:
\begin{align*}
  \hat{x}_k &= \mathrm{prox}_f(z_k), \\
  z_{k+1} &= z_k + \mathrm{prox}_g(2\hat{x}_k - z_k) - \hat{x}_k.
\end{align*}
\end{definition}

\begin{proposition}[Connection to ADMM]
ADMM applied to $\min f(x) + g(z)$ s.t.\ $x = z$ is equivalent to
Douglas-Rachford applied to the dual formulation.
\end{proposition}

\section{Applications in signal processing}

\begin{example}[LASSO -- sparse regression]
\index{LASSO}
The LASSO problem $\min_x \frac{1}{2}\norm{Ax-b}^2 + \lambda\norm{x}_1$
is composite with $f(x) = \frac{1}{2}\norm{Ax-b}^2$ (smooth, $L = \norm{A^\top A}$)
and $g(x) = \lambda\norm{x}_1$. ISTA gives:
\[
  x_{k+1} = S_{\lambda/L}\!\left(
    x_k - \frac{1}{L} A^\top(Ax_k - b)\right),
\]
where $S_\tau$ is component-wise soft thresholding.
\end{example}

\begin{example}[Total variation denoising]
For signal denoising of $y \in \R^n$:
\[
  \min_x \frac{1}{2}\norm{x - y}^2 + \lambda \norm{Dx}_1,
\]
where $D$ is the finite difference matrix. ADMM decomposes this into a
quadratic subproblem and soft thresholding.
\end{example}

\section{Exercises}

\begin{exercise}[$\star$ -- Proximal operators]
Compute the proximal operator of $g(x) = \lambda \norm{x}_1$ and of
$g(x) = \iota_{[0,+\infty)^n}(x)$ (indicator of the nonnegative orthant).
\end{exercise}

\begin{exercise}[$\star\star$ -- ISTA for LASSO]
Implement ISTA for LASSO with $A \in \R^{50 \times 200}$,
$x^* \in \R^{200}$ sparse (10 nonzero entries), and $b = Ax^* + \epsilon$.
Plot convergence and compare with FISTA.
\end{exercise}

\begin{exercise}[$\star\star$ -- Moreau identity]
Prove Moreau's identity: $\mathrm{prox}_g(v) + \mathrm{prox}_{g^*}(v) = v$.
Deduce $\mathrm{prox}_{\lambda\norm{\cdot}_\infty}$.
\end{exercise}

\begin{exercise}[$\star\star\star$ -- ADMM for LASSO]
Derive the ADMM iterations for LASSO. Show that the $x$-subproblem is a
linear system and the $z$-subproblem is soft thresholding. Implement and
compare with FISTA.
\end{exercise}

\begin{exercise}[$\star\star\star$ -- Accelerated convergence]
Show that the FISTA sequence $t_k$ satisfies $t_k \geq (k+1)/2$ and deduce
the $\mathcal{O}(1/k^2)$ convergence rate.
\end{exercise}
